\section{Non-tree curvature forces planar limits}
\label{sec:nontree-planar}

The geometric input in this section is a classical consequence of
Alexandrov angles and tangent cones.  Kleiner's dimension theorem gives
a broader account of Euclidean balls arising as rescaled subset limits
\cite[Theorem A(3), p.~2 of the author preprint]{Kleiner1999}.
We give the two-dimensional argument needed here, using the first
variation formula and the completed tangent-cone theorem of
\cite[II.3.6, II.3.18--3.19, pp.~185, 190--192]{BH1999}.
An $\mathbb R$-tree is a geodesic metric space in which every geodesic
triangle is a tripod.

\begin{lemma}\label{lem:nontree-angle}
If a $\cat$ space $X$ is not an $\mathbb R$-tree, some two nonconstant
geodesic segments with a common initial point have Alexandrov angle
strictly between $0$ and $\pi$.
\end{lemma}

\begin{proof}
Suppose instead that every such angle is $0$ or $\pi$.
Fix a unit-speed geodesic $\gamma:[0,L]\longrightarrow X$ and a point
$z$ outside its image.  The function
$f(t)=d(z,\gamma(t))$ is convex and $1$-Lipschitz.
For $0\le t<L$, the first variation formula gives
\[
 f'_+(t)=-\cos\angle_{\gamma(t)}(z,\gamma(L))\in\{-1,1\}.
\]
The right derivative of a convex function is nondecreasing.
It can therefore change from $-1$ to $1$ at most once.  Since $f$ is
Lipschitz and hence absolutely continuous, integration gives
\[
 f(t)=h+|t-a|\qquad(0\le t\le L)
\]
for some $a\in[0,L]$ and $h=f(a)\ge0$; the everywhere increasing and
everywhere decreasing cases correspond to $a=0$ and $a=L$.
Put $p=\gamma(a)$.  For every $t$ we have
\[
 d(z,\gamma(t))=d(z,p)+d(p,\gamma(t)).
\]
The uniqueness of geodesics in $X$ implies that $p$ lies on
$[z,\gamma(t)]$.  In particular the triangle with vertices
$z,\gamma(0),\gamma(L)$ is a tripod.  If $z$ lies on the original
segment, that conclusion is immediate.  Thus every geodesic triangle
is a tripod, proving the contrapositive.
\end{proof}

\begin{theorem}\label{thm:nontree-planar-limit}
Let $X$ be a $\cat$ space that is not an $\mathbb R$-tree.
For every finite tuple $z_1,\ldots,z_m\in\mathbb R^2$, there are
positive numbers $c_n$ and points $x_{1,n},\ldots,x_{m,n}\in X$ such
that
\begin{equation}\label{eq:nontree-planar-limit}
 \lim_{n\to\infty}c_n d_X(x_{i,n},x_{j,n})
   =\|z_i-z_j\|_2\qquad(1\le i,j\le m).
\end{equation}
No properness, local compactness, geodesic completeness, or finite
dimension is assumed.
\end{theorem}

\begin{proof}
By Lemma~\ref{lem:nontree-angle}, there are a point $p\in X$ and two
geodesic directions at $p$ whose angle is $\theta\in(0,\pi)$.
Let $\Sigma_pX$ be the completed space of directions.
The tangent-cone theorem says that $\Sigma_pX$ is $\mathrm{CAT}(1)$
and that its Euclidean cone is the completed $\cat$ tangent cone
$T_pX$ \cite[II.3.19, p.~191]{BH1999}.
The two directions are joined in $\Sigma_pX$ by a geodesic arc of
length $\theta$.  Distances on this arc agree with their parameter
differences.  Consequently its cone has distance formula
\[
 d\bigl((r,s),(r',s')\bigr)^2
 =r^2+(r')^2-2rr'\cos|s-s'|,
 \qquad s,s'\in[0,\theta],
\]
and is an isometrically embedded Euclidean sector of angle $\theta$.
Every finite planar set fits in this sector: a disk containing the
set can be translated sufficiently far along the sector's bisector.
We may therefore regard the given tuple as points
$v_1,\ldots,v_m\in T_pX$ with the same pairwise distances.

We next verify explicitly the finite approximation property of the
completed tangent cone.  Fix $\varepsilon>0$.  The cone of actual
geodesic germs is dense in $T_pX$, so choose germ vectors $w_i$ with
$d_{T_pX}(v_i,w_i)<\varepsilon$.
Represent each $w_i$ by a constant-speed geodesic
$\gamma_i:[0,t_0]\longrightarrow X$ starting at $p$, after shortening
their finitely many domains to one common interval.
The zero vector can be represented by the constant geodesic.
The definition of the Alexandrov angle, together with the Euclidean
comparison cosine law, gives
\[
 \lim_{t\downarrow0}\frac{d_X(\gamma_i(t),\gamma_j(t))}{t}
   =d_{T_pX}(w_i,w_j).
\]
Because there are only finitely many pairs, a single $t>0$ makes
each of these errors smaller than $\varepsilon$.
The triangle inequality in the tangent cone now gives
\[
 \left|\frac{d_X(\gamma_i(t),\gamma_j(t))}{t}
       -d_{T_pX}(v_i,v_j)\right|<3\varepsilon
 \qquad\text{for every }i,j.
\]
Apply this construction with $\varepsilon\downarrow0$, putting
$c_n=1/t_n$ and $x_{i,n}=\gamma_{i,n}(t_n)$.
This proves~\eqref{eq:nontree-planar-limit}.
The argument approximates only finitely many distance entries at a
time, so it does not require compactness of metric balls.
\end{proof}

\begin{corollary}\label{cor:closed-metric-planar-family}
Let $\mathcal M$ be a collection of finite pseudometrics closed under
restriction, positive rescaling, and pointwise limits on each fixed
finite label set.  If every finite subset of a non-tree $\cat$ space
has its metric in $\mathcal M$, then every finite Euclidean planar
metric belongs to $\mathcal M$.
\end{corollary}

\begin{proof}
Each rescaled finite distance table in
Theorem~\ref{thm:nontree-planar-limit} belongs to $\mathcal M$.
Its pointwise limit does as well.
\end{proof}

The conclusion concerns limits of rescaled finite configurations.
It does not assert that $X$ itself contains a flat square: a
hyperbolic plane, for example, is a non-tree $\cat$ space whose small
rescaled configurations approach Euclidean ones.  In the application
to prices, the remaining input is a $\cat$ realization whose entire
metric is a price on every finite subset.
